\unit{1962년 논평}
\sid{com}
\src{297}
\seminar{제14년차, 1961/62, 보충\\1962년 5월}
\paperhead{대수기하학의 기초.\\논평}%
  {알렉상드르 그로텐디크 지음}

\section*{일러두기}
\addcontentsline{toc}{section}{일러두기}

우리는 1957년부터 1962년까지 부르바키 세미나에서 대수기하학의 기초에
관한 여덟 강연을 하였다. 첫 번째를 제외한 이 강연들은 스킴의 틀 안에서
서술되어 있다. 여기서 진술한 모든 결과는 장 디외도네와 알렉상드르
그로텐디크의 『Éléments de géométrie algébrique』에 실릴 것이다. 그러나
이 강연들 가운데 어느 것의 내용도 현재 이미 출판되었거나 준비 중인 어느
장에도, 다른 어떤 책이나 논문에도 다루어져 있지 않으며, 아마 앞으로도
몇 해 동안은 그러할 것이다. 이것이 우리가 이 강연들을 한데 모은 주된
이유다. 이 강연들은 정식 집필본을 기다리는 동안 이용자들에게 스킴 이론의
몇몇 개념과 핵심 결과를 제공할 것이다. 더욱이 이 글들을 읽으면 체계적
논고에서 불가피하게 번잡해지는 세부에 방해받지 않고 그 결과와 개념에
빠르게 익숙해질 수 있으며, 그러한 논고를 공부하는 데 반드시 필요한
동기도 얻을 수 있다.

독자의 편의를 위해 강연별로 묶은 몇 가지 논평과 정정 사항을 여기에
모았다. 특히 이들은 각 글이 작성된 뒤 이루어진 진전을 곳곳에서 알리고
몇몇 추가 참고문헌을 제시한다.

이 강연들에 나오는 여러 결과는 IHES 대수기하학 세미나와 1961/62년에
하버드에서 연이어 열린 두 세미나에서 자세히 다루어졌다. 첫 번째 세미나는
내가, 두 번째 세미나는 멈퍼드--테이트가 진행했으며, 이 두 세미나의 등사판
강의록(리히텐바움 씨 작성)은 현재 준비 중이다.

\section*{사용한 약어}
\addcontentsline{toc}{section}{사용한 약어}

\begin{description}
\item[SGA :] 대수기하학 세미나, Institut des Hautes Études Scientifiques,
1960--1962.

\src{298}
\item[SHG :] 그로텐디크 세미나, 하버드 대학교, 1961/62.

\item[SHMT :] 멈퍼드--테이트 세미나, 하버드 대학교, 1961/62.

\item[EGA :] GROTHENDIECK (A.)와 DIEUDONNÉ (J.). --- Éléments de géométrie
algébrique. --- Paris, Presses universitaires de France, 1960, 1961, …
(Institut des Hautes Études Scientifiques, Publications mathématiques, 4, 8,
11, …).

\item[TDTE :] GROTHENDIECK (Alexander). --- Technique de descente et théorèmes
d'existence en géométrie algébrique, I--VI, \emph{Séminaire Bourbaki}, 제12권,
1959/60, 제190호와 제195호; 제13권, 1960/61, 제212호와 제221호; 제14권,
1961/62, 제232호와 제236호.
\end{description}

% FGA 정본 인쇄문 보존 장치: 원자적 통합 R1.
\par\smallskip
\begingroup\footnotesize
\noindent\textit{FGA-CANON-ERR-0030. 인쇄본의 독법은 C093에 계속 기록되어 있으며, FGA-CANON-DISP-0032가 현행 본문의 최소 수정을 규율한다.}\par
\endgroup
